\section{連立漸化式と固有値}\label{節：連立漸化式と固有値}
\subsection{和と差は新しい座標}

前節の行列で更新される二つの数列の一般項を求めよう.
\[
 \begin{aligned}
 x_{n+1}&=2x_n+y_n,\\
 y_{n+1}&=x_n+2y_n
 \end{aligned}
\]
を考える. 二式の和と差を取れば
\[
 \begin{aligned}
 x_{n+1}+y_{n+1}&=3(x_n+y_n),\\
 x_{n+1}-y_{n+1}&=x_n-y_n.
 \end{aligned}
\]
そこで$u_n=x_n+y_n,\,v_n=x_n-y_n$と置くと,\,二つの更新が独立する.
元へ戻す式は
\[
 x_n=\frac{u_n+v_n}{2},\qquad
 y_n=\frac{u_n-v_n}{2}.
\]
和と差を作ることは,\,情報を捨てることではなく,\,二つの値を別の座標で持つことだった.

\noindent\textbf{【例題】}\par\nobreak
上の連立漸化式を$x_1=1,\,y_1=0$から解け.

\noindent\textbf{【方針】}\par\nobreak
和は3倍,\,差は一定です. 二つを解いて元へ戻します.

\noindent\textbf{【解答】}\par\nobreak
\[
 x_n=\frac{3^{n-1}+1}{2},\qquad
 y_n=\frac{3^{n-1}-1}{2}.
\]
\noindent\textbf{【解法】}\par\nobreak
$u_1=v_1=1$なので$u_n=3^{n-1},\,v_n=1$である.
復元式に代入して一般項を得る.

行列$A=\begin{pmatrix}2&1\\1&2\end{pmatrix}$は,
$\begin{pmatrix}1\\1\end{pmatrix}$を3倍し,
$\begin{pmatrix}1\\-1\end{pmatrix}$を1倍する.
一般に零でない$\boldsymbol u$が
\[
 A\boldsymbol u=\lambda\boldsymbol u
\]
を満たすとき,$\boldsymbol u$を\textbf{固有ベクトル},
$\lambda$を対応する\textbf{固有値}という.
固有値は,その方向の一回の更新の倍率である.
負の倍率では向きが反転するが,原点を通る同じ直線上にとどまる.
零ベクトルはどんな$\lambda$でも式を満たすため,\,定義から除く.


\sectionref{節：座標と基底}の二方向が,\,更新に適した方向でもあった.
新しい座標は$\frac{u_n}{2},\frac{v_n}{2}$であり,一方は3倍,他方は1倍される.
和と差を選ぶ理由は,この二方向に分けると更新が混ざらなくなることにある.

二行二列の一般の行列
$A=\begin{pmatrix}p&q\\r&s\end{pmatrix}$
でも,$(\lambda I-A)\boldsymbol u=\boldsymbol0$に零でない解がある条件は
\[
 \begin{aligned}
 \det(\lambda I-A)
 &=(\lambda-p)(\lambda-s)-qr\\
 &=\lambda^2-(p+s)\lambda+ps-qr=0.
 \end{aligned}
\]
この方程式を行列$A$の\textbf{特性方程式}という.
行列式が0であることは,対応する連立一次方程式に零でない解があることを意味する.
前節の行列式を使えば,和と差を思い付かない場合も固有値を探せる.
ここで求めるのは,不動点の値ではなく各方向の倍率である.


\subsection{対角化は更新の分離}\label{節：対角化は更新の分離}

固有ベクトルを列に並べた
\[
 P=\begin{pmatrix}1&1\\1&-1\end{pmatrix}
\]
を使い,\,元の状態を$\boldsymbol z_n=P\boldsymbol c_n$と書く.
元の更新$\boldsymbol z_{n+1}=A\boldsymbol z_n$を代入すると
\[
 \boldsymbol c_{n+1}=P^{-1}AP\boldsymbol c_n.
\]
$P$で元へ戻し,$A$で更新し,$P^{-1}$で新しい座標に直す順序である.
この例では
\[
 P^{-1}AP=\begin{pmatrix}3&0\\0&1\end{pmatrix}=D.
\]
対角線の外の成分が零なので,\,新しい二成分が混ざらない.
このように座標を変えて対角行列にする操作を\textbf{対角化}という.

$AP=PD$から$A=PDP^{-1}$となり,
\[
 A^m=PD^mP^{-1}
\]
を得る. 隣り合う$P^{-1}P$が消えるためである.
$D^m$では対角成分を$m$乗するだけでよい.
公式は,\,\textbf{成分に分け,\,各々を更新し,\,元へ戻す}という解法そのものである.

一般の二成分でも,\,二つの異なる実固有値があれば実数の範囲で対角化できる.
異なる固有値の固有ベクトルが比例すると,\,同じ零でないベクトルに異なる倍率を指定することになり矛盾するからである.
固有値が非実数の場合は,\,実数の二方向には分けられない.
その場合は\chapterref{章：複素平面と漸化式}で平面の回転として扱う.

\subsection{多くの数列を同時に追う}

三つ以上の数列も同じようにまとめられる.
例えば
\[
 \begin{pmatrix}x_{n+1}\\y_{n+1}\\z_{n+1}\end{pmatrix}
 =
 \begin{pmatrix}2&1&1\\1&2&1\\1&1&2\end{pmatrix}
 \begin{pmatrix}x_n\\y_n\\z_n\end{pmatrix}
\]
では,\,全体の和$u_n=x_n+y_n+z_n$は4倍される.
一方,\,差$v_n=x_n-y_n$,$w_n=y_n-z_n$はそれぞれ変わらない.
三つの量から
\[
 \begin{aligned}
 x_n&=\frac{u_n+2v_n+w_n}{3},\\
 y_n&=\frac{u_n-v_n+w_n}{3},\\
 z_n&=\frac{u_n-v_n-2w_n}{3}
 \end{aligned}
\]
と元へ戻せるので,\,更新を分離できた.
初期値が$(1,0,0)$なら
\[
 x_n=\frac{4^{n-1}+2}{3},\qquad
 y_n=z_n=\frac{4^{n-1}-1}{3}.
\]
この行列では固有値1が重複するが,\,和が零の平面全体が変わらない.
その平面に独立な二方向を選び,4倍になる方向と合わせれば基底になる.
固有値の重複が,\,必ずしも分離を妨げるわけではない例である.

一般に$d$個の数列を$\boldsymbol z_n$にまとめれば,
$d$行$d$列の行列で$\boldsymbol z_{n+1}=A\boldsymbol z_n$と書ける.
$d$本の一次独立な固有ベクトルがあれば,\,それらを基底にして$d$個の等比型へ分けられる.
\chapterref{章：隣接多項間漸化式}の隣接多項間漸化式も,\,こうした連立更新の特別な場合に当たる.
